% ---------------------------------------------------------------------
\documentclass{article}

\usepackage{nberpreamble}
\title{\bfseries Plots Merge}
\author{\sffamily Prepared by Mauricio C\'aceres Bravo}
\date{\sffamily \today}

% ---------------------------------------------------------------------
\begin{document}
\displayoptions
\clearpage

% ---------------------------------------------------------------------
\section{What are we plotting?}
\label{sec:what_are_we_plotting_}

We aim to segment the data into 4 groups:
\begin{enumerate}
  \item Plan had above/below median alpha
  \item Plan was terminated or not-terminated
\end{enumerate}

Within a given year, each group assigned to a beneficiary will be
unique.  However, a beneficiary may change plans across years for
various reasons: If their plan terminates they would be forced into a
new plan, or if they move they would chanage plans by definition (since
we have determined that the "same" plan is different in different
counties).

Even if their plan does not terminate and they do not move, they may
appear with a new plan the next year. Hence this definition is not
unique across years. Since we want to compare a \textbf{\textit{cohort}}
of people, when we make the plots we fix a year, call it $t$, and we
segment the sample using plan alphas and terminations from year $t - 1$
using a single year of data. Hence we have a single data point per
beneficiary which indicates which of the 4 groups they belong to:
\begin{enumerate}[label=Group \arabic*:, leftmargin = 6em]
  \item Above median, non-terminated
  \item Above median, terminated
  \item Below median, non-terminated
  \item Below median, terminated
\end{enumerate}

We merge this single year of data into every other year, giving us a
way to track the cohort over time that is consistent. The resulting
plot will look like this:
\begin{figure}[!ht]
  \centering
  \caption{Average Plan Effects Over Time}
  \label{fig:average_plan_effects_over_time}
  \begin{tikzpicture}
    \begin{axis}[name=plot1
      %,title=
      ,width=15cm
      ,height=9cm
      ,ymin=-.6
      ,ymax=.8
      ,xmin=0
      ,xmax=2
      ,domain=0:2
      ,ylabel=Average Plan Effects
      ,y label style={
        at = {(axis description cs:-0.150,0.5)},
        rotate = 90
      }
      %,xlabel=$x$
      ,yticklabels = {-0.005,-0.003,-0.001,0.001,0.003,0.005}
      ,ytick       = {-0.5,-0.3,-0.1,0.1,0.3,0.5}
      ,xticklabels = {,,}
      ,xtick       = {0.05,1.95}
      ,every tick/.style = {
        black,
        semithick,
      }
    ]
    \draw [-, line width = 1.25pt, color=MaterialBlue300, dashed]
      (axis cs:0.05, 0.3) -- (axis cs:1.95, 0.2);
    \draw [-, line width = 1.25pt, color=MaterialBlue900]
      (axis cs:0.05, 0.275) -- (axis cs:1.95, -0.1);

    \draw [-, line width = 1.25pt, color=MaterialRed300, dashed]
      (axis cs:0.05, -0.5) -- (axis cs:1.95, -0.45);
    \draw [-, line width = 1.25pt, color=MaterialRed900]
      (axis cs:0.05, -0.55) -- (axis cs:1.95, -0.12);

    \node[above] at (axis cs:0.1, 0) {$t - 1$};
    \node[above] at (axis cs:1.95, 0) {$t$};

    % legend
    \draw [-, line width = 1.25pt, color=MaterialBlue300, dashed]
      (axis cs:1, 0.7) -- (axis cs:1.1, 0.7);
      \node[right] at (axis cs:1.1, 0.7) {Above median, non-terminated};
    \draw [-, line width = 1.25pt, color=MaterialBlue900]
      (axis cs:1, 0.6) -- (axis cs:1.1, 0.6);
      \node[right] at (axis cs:1.1, 0.6) {Above median, terminated};

    \draw [-, line width = 1.25pt, color=MaterialRed300, dashed]
      (axis cs:1, 0.5) -- (axis cs:1.1, 0.5);
      \node[right] at (axis cs:1.1, 0.5) {Below median, non-terminated};
    \draw [-, line width = 1.25pt, color=MaterialRed900]
      (axis cs:1, 0.4) -- (axis cs:1.1, 0.4);
      \node[right] at (axis cs:1.1, 0.4) {Below median, terminated};
    \end{axis}
  \end{tikzpicture}
\end{figure}

With each line representing the average plan effect in the indicated
year using the cohort groups for year $t$, which were defined using $t - 1$
variables.

% ---------------------------------------------------------------------
\section{Plots version 1 (``old'')}
\label{sec:plots_version_1_old_}

For a given beneficiary, we have the following:
\begin{table}[H]
  \centering
  \caption{Example of Data for Beneficiary 1}
  \label{tab:example_of_data_for_beneficiary_1}
  \begin{tabular}{cccccc}
  \toprule
  \verb`bene` & \verb`year` & \verb`plan` & \verb`ma` & \verb`term` & \verb`above median` \\
  \midrule
  1 & 2006 & 1 & 1 & 0 & 0 \\
  1 & 2007 & 1 & 1 & 0 & 1 \\
  1 & 2008 & 2 & 0 & 0 & 1 \\
  1 & 2009 & 2 & 0 & 0 & 0 \\
  1 & 2010 & 1 & 1 & 1 & 0 \\
  1 & 2011 & 2 & 0 & 0 & 1 \\
  \bottomrule
  \end{tabular}
\end{table}

For simplicity in this example, we write ``plan'' when we mean
county-plan, and we will consider a single year of data: 2011, which is
the year with the most terminations. Since 2011 is the reference year,
that is, $t = 2011$, the plan we care about is the $t - 1 = 2010$ plan,
which is plan \#1. Note that it does not matter that the 2011 plan is
no longer an MA plan, and that the beneficiary was not enrolled in an
MA plan in 2009. We only care about the $t - 1 = 2010$ data to assign
a group to this beneficiary, which in this case is ``Terminated, Below
Median'':  Group 4.

In general, for reference year 2011 we have 2010 data
\begin{table}[H]
  \centering
  \caption{Assigning Groups for Reference Year 2011}
  \label{tab:assigning_groups_for_reference_year_2011}
  \begin{tabular}{ccccccc}
  \toprule
  \verb`bene` & \verb`year` & \verb`plan` & \verb`ma` & \verb`term` & \verb`above median` & \verb`group` \\
  \midrule
  1 & 2010 & 1 & 1 & 1 & 0 & 4 \\
  2 & 2010 & 2 & 0 & 0 & . & . \\
  3 & 2010 & 3 & 1 & 1 & 0 & 4 \\
  4 & 2010 & 4 & 0 & 1 & . & . \\
  5 & 2010 & 5 & 1 & 1 & 0 & 4 \\
  6 & 2010 & 6 & 1 & 1 & 1 & 2 \\
  7 & 2010 & 7 & 1 & 0 & 0 & 3 \\
  8 & 2010 & 8 & 1 & 0 & 1 & 1 \\
  \bottomrule
  \end{tabular}
\end{table}

Note that non-MA plans are dropped, regardless of terminated status.
So they would neither be assigned an above/below median indicator nor
a group for the plots. To produce the $t - 1, t$ time series plots, we
would keep only \verb`bene` and \verb`group` from this data and merge by
\verb`bene` onto the main data, as illustrated in the following table
(note that group is time-invariant, even as other variables change):
\begin{table}[H]
  \centering
  \caption{Merge in Group Assignments for Reference Year 2011}
  \label{tab:merge_in_group_assignments_for_reference_year_2011}
  \begin{tabular}{ccccccc}
  \toprule
    \verb`bene` & \verb`year` & \verb`plan` & \verb`ma` & \verb`term` & \verb`above median` & \verb`group` \\
    \midrule
    1 & 2010 & 1 & 1 & 1 & 0 & 4 \\
    2 & 2010 & 2 & 0 & 0 & . & . \\
    3 & 2010 & 3 & 1 & 1 & 0 & 4 \\
    4 & 2010 & 4 & 0 & 1 & . & . \\
    5 & 2010 & 5 & 1 & 1 & 0 & 4 \\
    6 & 2010 & 6 & 1 & 1 & 1 & 2 \\
    7 & 2010 & 7 & 1 & 0 & 0 & 3 \\
    8 & 2010 & 8 & 1 & 0 & 1 & 1 \\
    \midrule
    1 & 2011 & 2 & 0 & 0 & . & 4 \\
    2 & 2011 & 2 & 0 & 0 & . & . \\
    3 & 2011 & 8 & 1 & 0 & 1 & 4 \\
    4 & 2011 & 8 & 1 & 0 & 1 & . \\
    5 & 2011 & 2 & 0 & 0 & . & 4 \\
    6 & 2011 & 7 & 1 & 0 & 0 & 2 \\
    7 & 2011 & 7 & 1 & 0 & 0 & 3 \\
    8 & 2011 & 8 & 1 & 0 & 1 & 1 \\
  \bottomrule
  \end{tabular}
\end{table}

% ---------------------------------------------------------------------
\section{Plots version 2 (current)}
\label{sec:plots_version_2_current_}

We always want to use $t - 1$ variables to compute the groups by
beneficiary. However, for the IV regressions, we use $t - 1$ variables
at year $t$: That is, we use lag variables in our main specification and
to construct our instrument. Hence we wanted to make a version of the
plots that was consistent with this paradigm. This is achievable by
replicating \Cref{tab:assigning_groups_for_reference_year_2011} using
$t - 1$ (lag) variables at year $t$:
\begin{table}[H]
  \centering
  \caption{Assigning Groups for Reference Year 2011}
  \label{tab:assigning_groups_for_reference_year_2011}
  \begin{tabular}{ccccccc}
  \toprule
  \verb`bene` & \verb`year` & \verb`lag plan` & \verb`lag ma` & \verb`lag term` & \verb`lag above median` & \verb`group` \\
  \midrule
  1 & 2011 & 1 & 1 & 1 & 0 & 4 \\
  2 & 2011 & 2 & 0 & 0 & . & . \\
  3 & 2011 & 3 & 1 & 1 & 0 & 4 \\
  4 & 2011 & 4 & 0 & 1 & . & . \\
  5 & 2011 & 5 & 1 & 1 & 0 & 4 \\
  6 & 2011 & 6 & 1 & 1 & 1 & 2 \\
  7 & 2011 & 7 & 1 & 0 & 0 & 3 \\
  8 & 2011 & 8 & 1 & 0 & 1 & 1 \\
  \bottomrule
  \end{tabular}
\end{table}

Note that the groupings are identical in this case, but the year
is now 2011 and we are instead using lag versions of the same
variables. From here it looks tempting to proceed just as in
\Cref{sec:plots_version_1_old_}: Keep \verb`bene` and \verb`group` then
merge by \verb`bene` onto the main data. However, we need to account
for deaths:  The above only defines the groups for \textit{surviving}
beneficiaries. We actually need to insert two additional steps:
\begin{enumerate}
  \item Rename \verb`lag plan` to \verb`plan` and replace \verb`year` with \verb`year - 1`.
  \item Merge by \verb`plan` (which has the \verb`lag plan` data) and year into the main data.
\end{enumerate}

This would assign each of the 4 groups in year $t - 1$ to all
beneficiaries enrolled in the above plans regardless of whether they
are dead in year $t$ or not. From here we \textit{can} proceed as in
\Cref{sec:plots_version_1_old_}, but it should be clear that this is
\textbf{\textit{not equivalent}} to version 1:
\begin{enumerate}

  \item If we define above/below median at year $t$ using plan
    effects from year $t - 1$, we are weighting by the number of
    \textit{surviving} beneficiaries from the plan, as opposed to
    by all beneficiaries who were enrolled.

  \item If every beneficiary in a terminated plan died, the plan would
    get dropped from the plots altogether.

  \item Since the regressions are meant to be ``stacked,'' that is,
    taken year-by-year, any selection criteria that is applied would
    be carried backwards, imposing a stronger selection criteria than
    intended for the plots. For example, if we requite \verb`lag ma` to
    equal 1, then we would exclude in year $t - 1$ every beneficiary who
    was not in an MA plan in year $t - 2$. Since we merge the groups
    by plan in year $t - 1$ to assign them to dead people, we are
    effectively imposing an ``always MA'' requirement as opposed to a
    ``lag MA'' restriction.

    Note that as long as we use the \verb`ifcommand()` option to
    apply lag restrictions (like \verb`lag ma == 1`) this last point
    should not be a problem.
\end{enumerate}

% ---------------------------------------------------------------------
\end{document}
